% Keep the temporary build directory out of the PDF trailer ID.
\ifdefined\pdfvariable\pdfvariable suppressoptionalinfo 512\relax\fi
\documentclass{ltjsbook}
\title{紙橋の型帳}
\author{紙橋模型工房}
\date{}
\begin{document}
\frontmatter
\maketitle
\tableofcontents
\mainmatter
\part{橋を組む}
\chapter{橋脚の型}
\section{折り筋}
細い帯を二つ折りにし、橋脚の芯を作る。\footnote{折り目は内側から押さえる。}
\subsection{採寸}
図\ref{fig:strip}の帯と表\ref{tab:cuts}の幅を照合する。
\begin{figure}[htbp]
\centering\rule{32mm}{3mm}
\caption{橋脚の帯}\label{fig:strip}
\end{figure}
\begin{itemize}
\item 白い帯には表の印を付ける。
\item 灰色の帯は補強に回す。
\end{itemize}
\newpage
\section{組立順}
\begin{enumerate}
\item 芯を台に立てる。
\item 桁を載せて傾きを見る。
\end{enumerate}
\chapter{試験台}
\section{間隔}
\subsection{切り幅}
式\ref{eq:span}で二本分の幅を出す。
\begin{equation}\label{eq:span}
s=2w
\end{equation}
\begin{table}[htbp]
\centering\caption{帯の切り幅}\label{tab:cuts}
\begin{tabular}{lr}部材 & 幅（ミリ）\\芯 & 6\\桁 & 8\end{tabular}
\end{table}
\appendix
\chapter{型の保管}
平らな封筒に帯を戻す。台の向きは型札に従う\cite{tag}。
\backmatter
\begin{thebibliography}{9}
\bibitem{tag} 紙橋模型工房『試験台の型札』。
\end{thebibliography}
\end{document}
